\subsection{Causal, constant-rate error cancellation}
\label{app:trace-proof}
\begin{proof}[Proof of \cref{thm:causal}]
\proofstep{(a) A valid error process.}
Write $S_i=\sum_{j=1}^i x_j$ and $V_i=N^2+N-i$ inside this proof.
For $N\le i\le2N$,
\[
 \frac{|x_i S_i|}{V_i}\le\frac{2N}{N^2-N}=\frac{2}{N-1}<\min\{2\eta,\gamma\}.
\]
Thus every conditional error probability in \cref{eq:causal-kernel} lies strictly between $0$ and $\tfrac{1}{2}$.
The Bayes predictor given the entire prefix and instruction is therefore $b_{i+1}x_i$ in both regimes.
The prescribed kernel depends only on the displayed past and current instruction, so the process is causal.
\proofstep{(b) Moments created by the error kernel.}
In the cancelling regime,
\begin{align*}
 \E[\xi_{i+1}\mid\bm x_{<i+1},b_{i+1}]&=\gamma-\frac{b_{i+1}x_i S_i}{V_i},\\
 \E[x_{i+1}\mid\bm x_{<i+1}]
 &=\frac{1}{2}\sum_{b=\pm1}b x_i\left(\gamma-\frac{b x_i S_i}{V_i}\right)
 =-\frac{S_i}{V_i}.
\end{align*}
At the seed, $\E[S_N^2]=N^2=V_N$. Inductively,
\begin{align*}
 \E[x_{i+1}S_i]&=-\frac{\E[S_i^2]}{V_i}=-1&&\text{(conditioning; induction)},\\
 \E[S_{i+1}^2]&=\E[S_i^2]+2\E[x_{i+1}S_i]+1
 =V_i-1=V_{i+1} &&\text{(definition of $V_i$)},\\
 \E[x_{i+1}S_{i+1}]&=\E[x_{i+1}S_i]+1=0&&\text{($x_{i+1}^2=1$)}.
\end{align*}
Global sign reversal preserves the kernel and reverses the displayed states.
Hence $\E[x_i]=\E[S_i]=0$ and each $x_i$ is uniform.
For $i\ge N+1$, the last identity gives
\[
 \E[S_i\mid x_i=1]=\E[S_i\mid x_i=-1]=0.
\]
Since the next instruction is independent of the prefix, averaging the error kernel conditional on $(x_i,b_{i+1})$ gives exactly $\eta$.
Also $\E[x_i(S_i-x_i)]=-1$, proving cancellation at every supervised position.
In the independent-error regime, a fresh uniform instruction makes $x_{i+1}$ uniform independently of its entire displayed past;
therefore $\E[x_i(S_i-x_i)]=0$ at the supervised positions.
Averaging either error kernel over its fresh instruction gives error probability $\eta$ at each of the $N+1$ noisy transitions.
The other $N-1$ transitions are correct; the expected incorrect fraction is $\tfrac{\eta(N+1)}{2N}$.
\proofstep{(c) High-probability corruption.}
Let $\mathcal F_i$ be the history generated by $(w,b_2,\xi_2,\ldots,b_i,\xi_i)$ immediately before drawing $b_{i+1}$,
with $\xi_2=\cdots=\xi_N=1$, and put $I_{i+1}=\mathbf 1\{\xi_{i+1}=-1\}$.
In either regime, averaging the fresh uniform instruction gives
\[
 \Pr(I_{i+1}=1\mid\mathcal F_i)
 =\frac{1}{2}\sum_{b=\pm1}\Pr(\xi_{i+1}=-1\mid\mathcal F_i,b_{i+1}=b)=\eta.
\]
For any $a_{N+1},\ldots,a_{2N+1}\in\{0,1\}$, repeated conditioning therefore yields
\[
 \Pr(I_{N+1}=a_{N+1},\ldots,I_{2N+1}=a_{2N+1})
 =\prod_{j=N+1}^{2N+1}\eta^{a_j}(1-\eta)^{1-a_j}.
\]
Thus $N_{\mathrm w}=\sum_{j=N+1}^{2N+1}I_j$ has the stated binomial law.
Write $M=\eta(N+1)$ only for this tail calculation. The lower tail satisfies
\begin{align*}
 \Pr\!\left(N_{\mathrm w}<\frac{\eta N}{2}\right)
 &\le \Pr\!\left(N_{\mathrm w}<\frac{M}{2}\right)&&\text{($\eta N\le M$)}\\
 &\le e^{M\log(2)/2}\E[2^{-N_{\mathrm w}}]&&\text{(Markov inequality)}\\
 &=e^{M\log(2)/2}\left(1-\frac{\eta}{2}\right)^{N+1}&&\text{(binomial law)}\\
 &\le\exp\!\left(-\frac{1-\log 2}{2}M\right)&&\text{($1-z\le e^{-z}$)}\\
 &\le\exp\!\left(-\frac{M}{8}\right)&&\text{($\log2\le\tfrac34$)}.
\end{align*}
Taking complements proves the high-probability bound.
The indicators are mutually independent; an error sign may still depend on its current instruction.
\proofstep{(d) Shared one-block training on all selected positions.}
For a prefix of length $i$, encode its source flag at position $i$ as in \cref{sec:architecture}.
The score coordinates and both MLP factors are shared across all $N$ prediction examples.
Write $\alpha_i(q)=\tfrac{e^q}{e^q+i-1}$. The matrix-derived attention read is
\[
 v_i=\alpha_i(q)x_i+\frac{1-\alpha_i(q)}{i-1}(S_i-x_i).
\]
Although the cancelling error sign depends on the instruction, averaging the latter gives
\[
 \E[x_{i+1}b_{i+1}v_i\mid\bm x_{<i+1}]=\gamma x_i v_i.
\]
The two averaged losses are consequently $-\tfrac{\gamma}{2}r_1r_2\bar \phi_\kappa(q)$, where
\begin{align}
 \bar \phi_1(q)&=\frac{1}{N}\sum_{i=N+1}^{2N}\frac{e^q}{e^q+i-1},\notag\\
 \bar \phi_0(q)&=\frac{1}{N}\sum_{i=N+1}^{2N}\frac{e^q-1}{e^q+i-1}.
 \label{eq:causal-flow-functions}
\end{align}
The initialization and chain rule from \cref{app:flow-proof} now give
\begin{align*}
 r_1=r_2=r>0,\qquad
 \ddt{r}&=\frac{\gamma}{2}r\bar \phi_\kappa(q),&
 \ddt{q}&=\gamma r^2\bar \phi_\kappa'(q),\\
 \bar F_\kappa(q)&:=\int_0^q\frac{\bar \phi_\kappa(v)}{\bar \phi_\kappa'(v)}\,dv,&
 r^2&=\varepsilon^2+\bar F_\kappa(q).
\end{align*}
Every source has attention at least $\tfrac{1}{2}$ exactly when $q\ge L:=\log(2N-1)$.
The fully correct control has the same independent-source moments, so its loss is $-\tfrac12r_1r_2\bar\phi_1(q)$.
Consequently its coordinate flow is exactly the independent-error flow with $\gamma=1$;
uniqueness gives $T_{\mathrm{independent}}=\tfrac{T_{\mathrm{clean}}}{\gamma}$.
The cancelling displayed-state law is also independent of $\eta$ on the admissible range:
its conditional mean in part (b), together with binary support, determines the same transition law for every $\eta$.
Its averaged loss changes only through $\gamma$, giving the same time rescaling within that range.
The hitting time is therefore
\[
 T=\int_0^L\frac{dq}{\gamma(\varepsilon^2+\bar F_\kappa(q))\bar \phi_\kappa'(q)}.
\]
\proofstep{(e) Uniform primitive and derivative bounds.}
For $0\le q\le L$ and $N+1\le i\le2N$,
\[
 \alpha_i(q)\ge\frac{e^q}{4N},\qquad \alpha_i(q)<\frac{2}{3},\qquad
 \bar \phi_\kappa'(q)\ge\frac{e^q}{12N}.
\]
For $0\le q\le1$, the derivative upper bounds are $\tfrac{e}{N}$ when $\kappa=1$ and $\tfrac{2e}{N}$ when $\kappa=0$.
A ratio of positive sums is a weighted average of the summand ratios. For $\kappa=1$ those ratios are
$\tfrac{1}{1-\alpha_i(q)}\in[1,3]$, so $q\le\bar F_1(q)\le3q$.
For $\kappa=0$ they are
\[
 (1-e^{-q})\frac{e^q+i-1}{i}.
\]
On $[0,1]$ these lie between $\tfrac{q}{2}$ and $2q$; throughout $[0,L]$ they are nonnegative and at most $3$.
Consequently
\[
 \frac{q^2}{4}\le\bar F_0(q)\le q^2\quad(0\le q\le1),\qquad
 \bar F_0(q)\ge\frac{1}{4}\quad(1\le q\le L),\qquad
 0\le\bar F_\kappa(q)\le3q.
\]
The prescribed gate scale gives $r^2\le \varepsilon^2+3L<B^2$ before routing, so the entire calculation is valid for the matrix network.
\proofstep{(f) Routing-time separation.}
The bounds in part (e) imply
\begin{align*}
 \frac{N}{3e\gamma}\log\!\left(1+\frac{3}{\varepsilon^2}\right)
 &\le T(\mathrm{independent})
 \le\frac{12N}{\gamma}\left[\log\!\left(1+\frac{1}{\varepsilon^2}\right)+\frac{1}{e}\right],\\
 \frac{N}{2e \gamma \varepsilon}\arctan\!\left(\frac{1}{\varepsilon}\right)
 &\le T(\mathrm{cancelling})
 \le\frac{12N}{\gamma}\left[\frac{\pi}{\varepsilon}+\frac{4}{e}\right].
\end{align*}
The upper estimates split the time integral at $q=1$ and use
$\int_0^\infty\tfrac{dq}{\varepsilon^2+\tfrac{q^2}{4}}=\tfrac{\pi}{\varepsilon}$ for the cancelled case.
Here $\int_0^\infty\tfrac{dv}{1+v^2}=\tfrac{\pi}{2}$.
The lower estimates integrate only over $[0,1]$.
Since $0<\varepsilon\le1$, these yield the asserted orders with absolute constants, including the clean order at $\gamma=1$.
For fixed admissible $N,\eta$,
\[
 \frac{T_{\mathrm{cancelling}}}{T_{\mathrm{independent}}}
 =\Theta\!\left(\frac{1}{\varepsilon\log(1+\varepsilon^{-2})}\right)\longrightarrow\infty.
\]
Together with $\gamma<1$, this proves the strict three-way ordering for sufficiently small initialization.
At $q=L$ the velocity remains positive and the gate inequality is strict, so the flow continues to $q>L$ on a nonempty interval.
There every source weight exceeds $\tfrac{1}{2}$, giving $\operatorname{sign}(z_i)=b_{i+1}x_i$ for all supervised prefixes.\qedhere
\end{proof}
